\cleardoublepage
\phantomsection
\addcontentsline{toc}{section}{Введение}
\section*{Введение}

Современные достижения в области анализа нейробиологических данных и нейровизуализации существенно расширили возможности разработки интерфейсов мозг--компьютер (ИМК)~\cite{wolpaw2000brain, makeig2004mining}. ИМК~--- это технология прямого взаимодействия нервной системы с внешними устройствами в обход периферических нервно-мышечных путей~\cite{mcfarland2002brain}. Регистрируя и декодируя нейробиологические сигналы~--- прежде всего электроэнцефалограмму (ЭЭГ), магнитоэнцефалограмму (МЭГ) и функциональную магнитно-резонансную томографию (фМРТ)~--- такие системы переводят ментальные намерения и когнитивные состояния в управляющие команды для протезов, роботов или компьютерных интерфейсов~\cite{song2020assistive, cruz2021self, dorin2025enhancing}. Для пациентов с тяжелыми двигательными или коммуникативными нарушениями это открывает возможность вернуть способность к автономному взаимодействию с окружающей средой~\cite{schwarz2020decoding}. Настоящая работа посвящена задаче декодирования и реконструкции визуальных стимулов по одновременным неинвазивным записям мозговой активности~--- фМРТ и ЭЭГ. Регистрируемые сигналы отражают динамику восприятия визуальных сцен, причем само восприятие зависит не только от свойств внешнего стимула, но и от внутренних когнитивных состояний, эмоциональных реакций и индивидуального опыта наблюдателя.

ЭЭГ и фМРТ~--- две ключевые нейробиологические модальности, чьи пространственно-временные характеристики разрешения взаимно дополняют друг друга. фМРТ измеряет связанный с нейронной активностью гемодинамический ответ через BOLD-сигнал, который проявляется с задержкой 4--8 секунд~\cite{dorin2024forecasting, Bandettini1992}; такая задержка отражает инерционность сосудистой системы, медленно реагирующей на возросшую потребность мозга в глюкозе~\cite{Ogawa1990, Logothetis2003}. Сильная сторона фМРТ~--- высокое пространственное разрешение, тогда как временное жестко ограничено малой частотой сканирования~\cite{Glover2011}. К этому добавляется чувствительность записи к артефактам от движений головы, пульсации, дыхания и теплового шума сканера~\cite{menon1999spatial, logothetis2008we}. ЭЭГ устроена противоположным образом: она дает высокое временное разрешение и позволяет следить за динамикой нейронных процессов в реальном времени~\cite{srinivasan1999methods}, однако ее пространственная локализация ограничена объемной проводимостью тканей головы, а сам сигнал ослаблен низким отношением сигнал/шум~\cite{signal_to_noise_eeg} и выраженной межиндивидуальной вариабельностью~\cite{morioka2015learning, jayaram2016transfer}. Совмещение ЭЭГ и фМРТ в одной сессии объединяет достоинства обеих модальностей, но заставляет считаться со спецификой магнитно-резонансной среды: сильное магнитное поле порождает мощные градиентные и кардиобаллистические артефакты~\cite{ritter2006simultaneous, niazy2005removal}. Для их компенсации необходимы магнитно-резонансно совместимые системы записи, точная синхронизация временных меток и специализированные методы постобработки~\cite{huster2012methods}.

В связи с развитием контрастивного обучения появилась возможность эффективно сопоставлять данные из разных модальностей. Так, в модели CLIP~\cite{radford2021learning} с помощью контрастивной функции потерь, построенной на косинусном расстоянии, удалось обучить совместное латентное пространство для эмбеддингов изображений и текстовых описаний. Аналогичный подход был распространен и на другие пары модальностей~--- при условии подбора соответствующих энкодеров~\cite{hu2023audio, wang2023connecting}. В области ИМК это привело к появлению методов реконструкции изображений по нейробиологическим сигналам, в первую очередь по фМРТ или ЭЭГ в отдельности~\cite{huang2021fmri, li2024visual}. Основная идея большинства современных моделей заключается в том, чтобы контрастивно приблизить эмбеддинги мозговой активности к CLIP-эмбеддингам изображений. Поскольку CLIP обучался на сопоставлении изображений и текста, его пространство эмбеддингов обладает богатой семантической структурой и демонстрирует высокую эффективность в задачах zero-shot (англ. zero-shot, <<без примеров>>), то есть на новые данные без дополнительного обучения~\cite{radford2021learning}. Это делает CLIP-эмбеддинги особенно подходящими для обуславливания генеративных моделей, используемых в реконструкции изображений, предъявленных испытуемому во время записи нейросигналов~\cite{li2024visual}. Однако существующие подходы опираются лишь на одну модальность. В то же время фМРТ и ЭЭГ демонстрируют фундаментально разные, но взаимодополняющие характеристики: первая обеспечивает высокое пространственное разрешение при низком временном, вторая~--- наоборот~\cite{Glover2011, srinivasan1999methods}. Их комбинация позволяет компенсировать слабые стороны каждой из модальностей и потенциально повысить качество реконструкции.

В данной работе решается задача реконструкции изображений по одновременно зарегистрированным сигналам фМРТ--ЭЭГ, как проиллюстрировано на рис.~\ref{fig:figure_0}.

\begin{figure}[h]
\centering
\includegraphics[width=1.\textwidth]{figure_0.pdf}
\caption{Иллюстрация задачи реконструкции изображений по одновременным сигналам фМРТ--ЭЭГ}
\label{fig:figure_0}
\end{figure}

До появления недавних мультимодальных наборов данных с визуальными стимулами, таких как~\cite{Telesford2023}, обработка ЭЭГ и фМРТ, как правило, проводилась раздельно. Новые датасеты открывают возможность интегрировать оба типа сигналов в единую архитектуру. По аналогии с предшествующими работами, в настоящей работе используется двухэтапный метод декодирования мозговой активности: 1)~контрастивное выравнивание сигналов фМРТ и ЭЭГ с эмбеддингами изображений; 2)~обуславливание предобученной модели генерации изображений на уточненный априорной моделью эмбеддинг сигналов мозга. Решение строится в три последовательных этапа, отражающих постепенное усложнение учитываемой структуры сигнала: анализ временной динамики фМРТ, последующий учет его пространственно-временной структуры в условиях ограниченного объема выборки и мультимодальное декодирование по совместным записям фМРТ--ЭЭГ. Основной же вклад настоящей работы заключается в следующем:
\begin{itemize}
    \item Задача декодирования визуальной информации по сигналам мозга формализована в единой постановке, объединяющей три подзадачи: оценку гемодинамической задержки, классификацию временных рядов фМРТ в условиях ограниченного объема выборки и мультимодальную реконструкцию стимулов по сигналам фМРТ--ЭЭГ.
    \item Предложена линейная модель прогнозирования фМРТ-отклика по визуальному стимулу, обеспечивающая оценку гемодинамической задержки BOLD-сигнала; полученная оценка согласуется с известным физиологическим диапазоном и используется для временной синхронизации модальностей.
    \item Разработан метод пространственно-временной классификации временных рядов фМРТ, основанный на масках активности и римановой геометрии ковариационных матриц и сохраняющий устойчивость в условиях ограниченного объема выборки одного испытуемого.
    \item Вводится мультимодальная модель, использующая одновременно и фМРТ, и ЭЭГ сигналы для реконструкции изображений, просматриваемых пациентом.
    \item Предлагается архитектура мультимодальной модели, в которой обучаются три модуля: два энкодера, независимо обрабатывающих наблюдаемые сигналы, и модуль объединения, переводящий их в единое комбинированное представление.
    \item Проводится сравнительный анализ исключением модальностей, демонстрирующий необходимость учета обоих сигналов фМРТ и ЭЭГ совместно.
\end{itemize}

\paragraph{Публикации и апробация результатов работы.}
По теме исследования автором подготовлены и представлены следующие работы.

\textit{Публикации в научных журналах:}
\begin{enumerate}
    \item \textit{\textbf{Dorin D.}, Kiselev N. et al.} Forecasting fMRI Images From Video Sequences: Linear Model Analysis // Health Information Science and Systems.~--- 2024. doi: \href{https://doi.org/10.1007/s13755-024-00315-5}{10.1007/s13755-024-00315-5}
    \item \textit{\textbf{Dorin D.}, Grabovoy A., Strijov V.} Enhancing fMRI Data Decoding with Spatiotemporal Characteristics in Limited Dataset // Doklady Mathematics.~--- 2025. doi: \href{https://doi.org/10.1134/S1064562425700383}{10.1134/S1064562425700383}
    \item \textit{\textbf{Dorin~D.}, Kiselev~N., Grabovoy~A.} Decoding Visual Information from Neural Signals: Image Reconstruction Based on Joint fMRI and EEG Analysis // Informatika i ee Primeneniya.~--- 2026.
\end{enumerate}

\textit{Доклады на научных конференциях:}
\begin{enumerate}
    \item \textit{\textbf{Дорин~Д.\,Д.}, Киселев Н.\,С., Грабовой А.\,В.} Пространственно-временные методы анализа временных рядов~// 66-я Всероссийская научная конференция МФТИ~--- 2024.

    \item \textit{\textbf{Дорин~Д.\,Д.}, Грабовой А.\,В.} Улучшение декодирования фМРТ в условиях ограниченной выборки~// 67-я Всероссийская научная конференция МФТИ~--- 2025.
    \item \textit{\textbf{Дорин~Д.\,Д.}, Грабовой~А.\,В.} Декодирование визуальной информации из сигналов мозга на основе пространственно-временных характеристик~// 68-я Всероссийская научная конференция МФТИ~--- 2026.
    \item \textit{\textbf{Дорин~Д.\,Д.}, Киселев~Н.\,С., Грабовой~А.\,В.} Декодирование визуальных стимулов из мультимодальных сигналов мозга~// 22-я конференция <<Математические методы распознавания образов>> (ММРО-2025).
    \item \textit{\textbf{Дорин Д.\,Д.}, Грабовой А.\,В.} Улучшение декодирования данных фМРТ в условиях ограниченного набора данных // AI Journey 2025.
\end{enumerate}
